\documentclass[12 pt, letterpaper]{hmcpset}
\usepackage{hw-preamble}
\setlist[enumerate, 1]{label = (\roman*)}

\name{Jayden Thadani}
\class{MATH 145 --- Topics in Geometry and Topology}
\assignment{Homework 1}
\duedate{29th January 2025}

\begin{document}

\begin{problem}[1.]
	The following picture is my attempt to draw a $1$-cycle $\alpha$, shown in blue.
	\begin{enumerate}
		\item In fact, I have made mistakes, Specifically, there are exactly three arrows pointing in the wrong direction. Can you identify with certainty which three arrows have been reversed?
		\item Having corrected the errors, compute $w(\alpha, \mathbf{x})$ on each region of $\mathbb{R}^{2} \setminus \lvert \alpha \rvert$.
	\end{enumerate}
\end{problem}

\pagebreak

\begin{problem}[2.]
	\begin{enumerate}
		\item Let $\gamma = \sum_{i = 1}^{n}$ be a $1$-cycle, and let $\mathbf{x} \in \mathbb{R}^{2} \setminus \lvert \gamma \rvert$. Show that the inequality
			\[
				\sum_{i = 1}^{n} \lvert \theta([\mathbf{a}_{i}, \mathbf{b}_{i}], \mathbf{x}) \rvert < 2 \pi
			\]
			implies that $w(\gamma, \mathbf{x}) = 0$.
		\item Let $\mathbf{a}, \mathbf{b}$ be points that lie in a circular disc of radius $r$ and let $\mathbf{x}$ be a point at distance $R > r$ from the centre of the disc. Using a simple trigonometric argument, show:
			\[
				\lvert \theta([\mathbf{a}, \mathbf{b}], \mathbf{x}) \rvert \le 2 \arcsin(r / R)
			\]
		\item Let $\gamma = \sum_{i = 1}^{n} [\mathbf{a}_{i}, \mathbf{b}_{i}]$ be a $1$-cycle such that all the points $\mathbf{a}_{i}, \mathbf{b}_{i}$ are contained in a circular disc of radius $r$. Let $\mathbf{x}$ be a point at distance $R > r$ from the centre of the disc. Suppose
			\[
				R > \frac{r}{\sin{(\pi / n)}}.
			\]
			Combine your results from (i) and (ii) to show that $w(\gamma, \mathbf{x}) = 0$.
		\item Explain why a stronger result than (iii) is totally obvious from the ray-escape formula.
	\end{enumerate}
\end{problem}

\pagebreak

\begin{problem}[3.]
	Here are a few of my favourite $1$-cycles. Let $n$ be a positive integer. Define points $v_{\ell}$ on the unit circle by the following formula:
	\[
		v_{\ell} = e^{2 \pi i \ell / n}
	\]
	Note that these are periodic in that $v_{\ell + n} = v_{\ell}$ for all $\ell$. Let $k$ be any integer. Define the \emph{regular cycle} $E(n, k)$ as follows:
	\[
		E(n, k) = [v_{1}, v_{1 + k}] + [v_{2}, v_{2 + k}] + \dots + [v_{n}, v_{n + k}].
	\]
	In each of the following cases, draw a sketch of the $1$-cycle and write the winding numbers in each component of the complementary region.
	\begin{enumerate}
		\item $E(5, 1)$ 
		\item $E(5, 2)$ 
		\item $E(5, 3)$ 
		\item $E(5, 3)$ 
		\item $E(6, 3)$ 
		\item $E(8, 2)$ 
		\item $E(8, 3)$ 
		\item $E(2, 1)$
		\item $E(5, -2)$.
		\item Repeat for one more $E(n, k)$ of your choice.
	\end{enumerate}
\end{problem}

\pagebreak

\begin{problem}[4.]
	Obtain a simple general formula for $w(E(n, k), 0)$, under the assumption that $0 < k < n / 2$.
\end{problem}

\pagebreak

\begin{problem}[5.]
	Consider the following question: what is the maximum possible value of $w(\gamma, \mathbf{x})$, when $\gamma$ is a $1$-cycle with $n$ edges and $\mathbf{x} \notin \lvert \gamma \rvert$? Let's call that number $w_{\text{max}}(n)$.
	\begin{enumerate}
		\item Let $n$ be a positive integer. Show that $\lvert w(\gamma, \mathbf{x}) \rvert < n / 2$, for all $1$-cycles $\gamma$ with $n$ edges.
		\item For all $n$, exhibit a $1$-cycle $\gamma$ and a point $\mathbf{x} \in \mathbb{R}^{2} \setminus \lvert \gamma \rvert$ such that $w(\gamma, \mathbf{x}) = \lfloor \frac{n - 1}{2} \rfloor$.
	\end{enumerate}
\end{problem}

\pagebreak

\begin{problem}[6.]
	Let $\gamma$ be a $1$-cycle in $\mathbb{R}^{2}$. We define the \emph{area} of $\gamma$ as follows
	\[
		\operatorname{area}(\gamma) = \int_{\mathbb{R}^{2}} w(\gamma, \mathbf{x}) \, d\mathbf{x}.
	\]
	Equivalently, partition the open set $\mathbb{R}^{2} \setminus \lvert \gamma \rvert$ into its connected components $U_{j}$, and let $w_{j}$ be the value of $w(\gamma, \mathbf{x})$ for $\mathbf{x} \in U_{j}$. Then
	\[
		\operatorname{area}(\gamma) = \sum_{j} w_{j} \operatorname{area}(U_{j}) 
	\]
	where $\operatorname{area}(U_{j})$ is simply the ordinary area of the region $U_{j}$.
	\begin{enumerate}
		\item Compute $\operatorname{area}(\alpha)$, for the (corrected) $1$-cycle $\alpha$ given in question 1.
		\item Construct a quadrilateral $\kappa = [\mathbf{a}, \mathbf{b}] + [\mathbf{b}, \mathbf{c}] + [\mathbf{c}, \mathbf{d}] + [\mathbf{d}, \mathbf{a}]$ with $\operatorname{area}(\kappa) = 0$, such that the four points $\mathbf{a}, \mathbf{b}, \mathbf{c}, \mathbf{d}$ are distinct and do not all lie on a straight line.
	\end{enumerate}
\end{problem}

\pagebreak

\begin{problem}[7.]
	Give a 1-line proof that $\operatorname{area}(\gamma_{1} + \gamma_{2}) = \operatorname{area}(\gamma_{1}) + \operatorname{area}(\gamma_{2})$ for any two $1$-cycles $\gamma_{1}, \gamma_{2}$.
\end{problem}

\pagebreak

\begin{problem}[8.]
	\begin{enumerate}
		\item Let $\gamma = \sum_{i = 1}^{n} [\mathbf{a}_{i}, \mathbf{b}_{i}]$ be a $1$-chain, and let $\mathbf{y} \in \mathbb{R}^{2}$. Show that the $1$-chain
			\[
				\hat{\gamma} = \sum ([\mathbf{a}_{i}, \mathbf{b}_{i}] + [\mathbf{b}_{i}, \mathbf{y}] + [\mathbf{y}, \mathbf{a}_{i}])
			\]
			is always a $1$-cycle.
		\item Suppose $\gamma$ itself is a $1$-cycle. Show that $w(\gamma, \mathbf{x}) = w(\hat{\gamma}, \mathbf{x})$ for all $\mathbf{x} \in \mathbb{R}^{2} \setminus \lvert \hat{\gamma} \rvert$.
		\item Deduce that
			\[
				\operatorname{area}(\gamma) = \frac{1}{2} \sum_{i = 1}^{n} \lvert \mathbf{y} - \mathbf{a}_{i} \rvert \cdot \lvert \mathbf{y} - \mathbf{b}_{i} \rvert \sin{(\theta([\mathbf{a}_{i}, \mathbf{b}_{i}], \mathbf{y}))}.
			\]
	\end{enumerate}
\end{problem}

\pagebreak

\begin{problem}[9.]
	Calculate $\operatorname{area}(E(n, k))$.
\end{problem}

\end{document}
